\section{Introduction}
\label{sec:introduction}

Scientific agents—systems that propose, execute, and evaluate computational
experiments—produce trajectories that are substantially more complex than the
single-turn responses studied in standard language-model evaluation.
A typical long-horizon scientific run involves dozens of tool invocations,
code edits, data transformations, metric computations, and evidence-gathering
steps before arriving at a final scientific claim. Yet the dominant evaluation
paradigm asks only whether that terminal claim is correct. This practice
masks critical questions: when did the trajectory begin to diverge from a
valid scientific procedure? Which causal mechanism triggered the divergence?
Could the failure have been detected early enough to intervene? And does
recovery remain possible once a warning has been issued?

The agent-reliability community has produced valuable resources—failure
taxonomies~\citep{mast2025,aegis2025}, trajectory-diagnosis
tools~\citep{agentrx2026}, scientific-agent benchmarks~\citep{mlrbench2025,sciagentarena2026},
and execution-feedback training methods~\citep{capf2026,afterburner2025}.
These contributions establish that agents fail, that failures have structure,
and that trajectories contain diagnostic information. However, no existing
resource combines governed causal fault injection, structured
state-difference observation, independent scientific adjudication, and
frozen transfer splits into a single measurement artifact that can answer
the four questions above with reproducible, uncertainty-quantified evidence.

This paper makes three contributions. First, we define a governed trace
contract and causal fault model with seven mechanism families, deterministic
injectors, clean counterfactuals, and separate execution-versus-scientific
adjudication—a measurement infrastructure that can be reused across models,
tasks, and fault regimes. Second, we construct a benchmark spanning four
workload levels and six held-out transfer splits, release it with documented
access, licensing, provenance, and maintenance plans, and report failure
hazard, propagation, and recoverability measurements derived from controlled
injection. Third, we build and evaluate a calibrated early-warning and
attribution baseline that predicts terminal scientific failure from
execution-grounded state-difference histories and evidence provenance,
tested on held-out task families, unseen fault injector implementations,
and unseen fault compositions.

The paper is deliberately a benchmark and measurement contribution. The
warning model uses a compact, interpretable baseline; learned recovery
policies and process-supervised post-training belong to a companion method
paper that consumes the frozen artifact described here. The benchmark is
structured so that negative results—a failure law that does not replicate,
a warning signal that does not transfer—are recorded and released rather
than silently discarded.
